\documentclass[11pt]{article}
\usepackage[a4paper,margin=24mm]{geometry}
\usepackage{fontspec}
\usepackage{xcolor}
\usepackage{amsmath,amssymb}
\usepackage[hidelinks]{hyperref}
\setmainfont[Path=fonts/,Extension=.otf,UprightFont=*-Regular,BoldFont=*-Bold]{NotoSerifCJKkr}
\definecolor{reviewercolor}{RGB}{38,82,126}
\definecolor{responsecolor}{RGB}{36,113,78}
\definecolor{revisioncolor}{RGB}{145,82,24}
\setlength{\parindent}{0pt}
\setlength{\parskip}{0.55em}
\setlength{\fboxsep}{7pt}
\setlength{\emergencystretch}{3em}
\renewcommand{\familydefault}{\rmdefault}
\newcommand{\fieldheading}[2]{\par\medskip\noindent{\color{#1}\rule{3pt}{1.25em}}\hspace{0.55em}\textbf{#2}\par\smallskip}
\newenvironment{reviewitem}[1]{\subsection*{#1}}{\par\medskip\hrule\bigskip}
\newenvironment{reviewercomment}{\fieldheading{reviewercolor}{1. 리뷰어 코멘트}\begin{quote}}{\end{quote}}
\newenvironment{authorresponse}{\fieldheading{responsecolor}{2. 저자 답변}}{\par}
\newenvironment{manuscriptrevision}[1]{\fieldheading{revisioncolor}{3. 논문 반영 결과}\textbf{수정 위치:} #1\par\smallskip}{\par}
\newcommand{\manuscripttitle}{연속 인과관계(CoC) 주석의 시간·조건 일관성 검사 프레임워크: 적용 가능성과 근거 추적}
\newcommand{\manuscriptid}{확인 필요 (제공 기록에서 확인되지 않음)}
\newcommand{\revisionround}{확인 필요 (개정 소스 버전: revision-v01)}
\newcommand{\responsedate}{2026년 10월 1일}
% Generated from the final manuscript auxiliary labels.
\expandafter\def\csname mloc@sec:intro\endcsname{1절(1쪽)}
\expandafter\def\csname mloc@sec:related\endcsname{2절(1쪽)}
\expandafter\def\csname mloc@fig:overview-pipeline\endcsname{그림 1(2쪽)}
\expandafter\def\csname mloc@sec:data\endcsname{3절(2쪽)}
\expandafter\def\csname mloc@sec:provenance\endcsname{3.1절(2쪽)}
\expandafter\def\csname mloc@sec:selection\endcsname{3.2절(2쪽)}
\expandafter\def\csname mloc@sec:artifacts\endcsname{3.3절(2쪽)}
\expandafter\def\csname mloc@sec:problem\endcsname{4절(2쪽)}
\expandafter\def\csname mloc@tab:related-position\endcsname{표 1(3쪽)}
\expandafter\def\csname mloc@tab:provenance\endcsname{표 2(3쪽)}
\expandafter\def\csname mloc@sec:error-types\endcsname{4.1절(3쪽)}
\expandafter\def\csname mloc@eq:hold-conflict\endcsname{식 1(3쪽)}
\expandafter\def\csname mloc@eq:missing-evidence\endcsname{식 5(3쪽)}
\expandafter\def\csname mloc@sec:semantics\endcsname{4.2절(3쪽)}
\expandafter\def\csname mloc@eq:active-update\endcsname{식 7(3쪽)}
\expandafter\def\csname mloc@sec:method\endcsname{5절(3쪽)}
\expandafter\def\csname mloc@sec:extraction\endcsname{5.1절(3쪽)}
\expandafter\def\csname mloc@sec:models\endcsname{5.2절(4쪽)}
\expandafter\def\csname mloc@lst:condition-step\endcsname{목록 1(4쪽)}
\expandafter\def\csname mloc@sec:queries\endcsname{5.3절(4쪽)}
\expandafter\def\csname mloc@lst:queries\endcsname{목록 2(4쪽)}
\expandafter\def\csname mloc@sec:experiments\endcsname{6절(4쪽)}
\expandafter\def\csname mloc@sec:legacy-test\endcsname{6.1절(4쪽)}
\expandafter\def\csname mloc@sec:scenario-design\endcsname{6.2절(4쪽)}
\expandafter\def\csname mloc@tab:worked-example\endcsname{표 3(5쪽)}
\expandafter\def\csname mloc@tab:transition\endcsname{표 4(5쪽)}
\expandafter\def\csname mloc@tab:scenario-variants\endcsname{표 5(5쪽)}
\expandafter\def\csname mloc@sec:retest-design\endcsname{6.3절(5쪽)}
\expandafter\def\csname mloc@sec:aux-design\endcsname{6.4절(5쪽)}
\expandafter\def\csname mloc@tab:regression\endcsname{표 6(5쪽)}
\expandafter\def\csname mloc@sec:results\endcsname{7절(5쪽)}
\expandafter\def\csname mloc@sec:regression-results\endcsname{7.1절(5쪽)}
\expandafter\def\csname mloc@sec:first-results\endcsname{7.2절(5쪽)}
\expandafter\def\csname mloc@sec:retest-results\endcsname{7.3절(5쪽)}
\expandafter\def\csname mloc@tab:first-results\endcsname{표 7(6쪽)}
\expandafter\def\csname mloc@tab:failures\endcsname{표 8(6쪽)}
\expandafter\def\csname mloc@tab:retest\endcsname{표 9(6쪽)}
\expandafter\def\csname mloc@tab:implementation\endcsname{표 10(6쪽)}
\expandafter\def\csname mloc@tab:agreement\endcsname{표 11(6쪽)}
\expandafter\def\csname mloc@sec:agreement-results\endcsname{7.4절(6쪽)}
\expandafter\def\csname mloc@sec:human-results\endcsname{7.5절(6쪽)}
\expandafter\def\csname mloc@tab:reviewer-distribution\endcsname{표 12(6쪽)}
\expandafter\def\csname mloc@sec:trajectory-results\endcsname{7.6절(6쪽)}
\expandafter\def\csname mloc@sec:discussion\endcsname{8절(6쪽)}
\expandafter\def\csname mloc@tab:disagreement-examples\endcsname{표 13(7쪽)}
\expandafter\def\csname mloc@tab:trajectory-cases\endcsname{표 14(7쪽)}
\expandafter\def\csname mloc@sec:threats\endcsname{9절(7쪽)}
\expandafter\def\csname mloc@sec:conclusion\endcsname{10절(8쪽)}

\newcommand{\mloc}[1]{\csname mloc@#1\endcsname}
\newcommand{\responsefront}[2]{%
\begin{center}{\LARGE\bfseries 심사 의견 답변서\par}\vspace{1em}{\large\bfseries\manuscripttitle\par}\end{center}
\vspace{0.8em}
\begin{tabular}{@{}ll@{}}논문 접수번호: & \manuscriptid \\ 수정 차수: & \revisionround \\ 답변일: & \responsedate\end{tabular}
\section*{답변에 앞서}
편집위원장님과 심사위원님께,

연속 CoC의 검사 범위, 의미 변환과 평가 근거를 구체화할 수 있도록 의견을 주셔서 감사합니다.
이번 개정은 자연 오류 성능 대신 시간·조건 관계의 검사 적용 가능성과 근거 추적으로 주장을 제한했습니다.
이미 수행한 재분석과 통제 시나리오의 최초 결과, 보완 후 동일 사례 재시험을 구분하여 반영했습니다.
사람 재평가나 개발 미사용 독립 시험을 수행했다고 주장하지 않습니다.
개정 원고의 파란색은 원 저장소의 최초 제출본 대비 이번에 추가·수정한 내용을 나타냅니다.
검정색 clean 원고는 동일 내용입니다. 아래 절·쪽·표 번호는 이번 개정 원고를 기준으로 합니다.
접수번호와 실제 심사 수정 차수는 기록이 없어 확인 필요로 남겼습니다.

\section*{심사위원 #1}
\textbf{총평(요지)}\par
\begin{quote}#2\end{quote}}


\begin{document}

\responsefront{2}{서론의 문장과 구조, 전반적인 국문 표현, UPPAAL 명칭, 연구 질문과 기여의 관계 및 관련 연구 대비 기술적 차이를 명확히 할 필요가 있다. 원문에서 반복된 의견은 한 번만 수록하고 문장 수정1.1·1.2는 각각 답변한다.}

\clearpage

\begin{reviewitem}{R2-1.1. CoC 간 양립 문장}

\begin{reviewercomment}

기존 문장: ``문제는 같은 장면의 CoC들이 개별적으로는 자연스러워도 시간순으로 함께 읽으면 양립하지 않을 수 있다는 점이다.'' 수정 제안: ``하지만 이 방식은 같은 장면의 CoC들이 개별적으로는 자연스러워도 시간순으로 함께 읽으면 양립하지 않을 수 있다는 문제점이 있다.

\end{reviewercomment}

\begin{authorresponse}

제안 문장을 서론에 반영했습니다. 이어 생성·작성 과정에서 조건 누락과 앞선 요구에 충돌하는 행동이 생길 수 있다는 이유를 설명했습니다. 다음 문단의 보행자 정지→녹색 신호 출발은 자연 자료에서 발견한 사례가 아니라 동기를 위한 가상 예시로 명시했습니다.

\end{authorresponse}

\begin{manuscriptrevision}{\mloc{sec:intro}; 아래 발췌: 1쪽}

본문에 다음 내용을 반영했습니다.

\begin{quote}

하지만 이 방식은 같은 장면의 CoC들이 개별적으로는 자연스러워도 시간순으로 함께 읽으면 양립하지 않을 수 있다는 문제점이 있다.

\end{quote}

\end{manuscriptrevision}

\end{reviewitem}

\clearpage

\begin{reviewitem}{R2-1.2. 불일치의 영향 문장}

\begin{reviewercomment}

기존 문장: ``이러한 불일치는 학습 자료의 행동 요구를 모호하게 하고 같은 자료를 다시 검사했을 때 일관된 결과를 얻기 어렵게 한다.'' 수정 제안: ``이러한 불일치 때문에 학습 자료의 행동 요구가 모호해지고, 같은 자료를 다시 검사했을 때 일관된 결과를 얻기 어렵다.

\end{reviewercomment}

\begin{authorresponse}

제안 표현으로 수정했습니다. 행동 요구의 모호성이 검토 동기임을 설명하되 이번 검사로 실제 학습 성능이 향상됐다고 주장하지 않습니다. 명시된 의무·증거의 일관성과 차량 안전 및 학습 효과를 구분했습니다.

\end{authorresponse}

\begin{manuscriptrevision}{\mloc{sec:intro}; \mloc{sec:conclusion}; 아래 발췌: 1쪽}

본문에 다음 내용을 반영했습니다.

\begin{quote}

이러한 불일치 때문에 학습 자료의 행동 요구가 모호해지고, 같은 자료를 다시 검사했을 때 일관된 결과를 얻기 어렵다.

\end{quote}

\end{manuscriptrevision}

\end{reviewitem}

\clearpage

\begin{reviewitem}{R2-2. 논문 전반의 국문 표현}

\begin{reviewercomment}

논문 전체에 걸쳐 영문을 번역한 듯한 문장이 다수 있으므로, 자연스러운 국문 문장으로 수정할 필요가 있다.

\end{reviewercomment}

\begin{authorresponse}

서론뿐 아니라 관련 연구·자료·문제 정의·방법·평가·결과·논의·한계·결론을 새 근거에 맞춰 다시 교열했습니다. 주어와 판단 대상을 명시하고 ``해제 근거 부족'', ``해당 의무와 양립'', ``개발 후 동일 사례 재시험''의 용어를 일관되게 사용했습니다. 자동 문구 추출과 작성 증거, 주석 문제와 처리 오류를 한 문장에 섞지 않도록 정리했습니다.

독립 검증이나 합의 과정처럼 수행되지 않은 작업을 완료형으로 표현하지 않았습니다. 영문 초록과 결과 표의 분모도 국문 본문과 맞췄습니다.

\end{authorresponse}

\begin{manuscriptrevision}{\mloc{sec:data}; \mloc{sec:problem}; \mloc{sec:method}; \mloc{sec:results}; \mloc{sec:discussion}; \mloc{sec:threats}; \mloc{sec:conclusion}; 아래 발췌: 3쪽}

본문에 다음 내용을 반영했습니다.

\begin{quote}

정지나 대기 중 해제 여부가 미상이라는 이유만으로 문제를 기록하지 않는다. 진행할 때 남은 의무가 거짓이면 모순, 미상이면 근거 부족이다.

\end{quote}

\end{manuscriptrevision}

\end{reviewitem}

\clearpage

\begin{reviewitem}{R2-3. UPPAAL 명칭}

\begin{reviewercomment}

UPPAAL은 무엇의 약자인가?

\end{reviewercomment}

\begin{authorresponse}

UPPAAL은 Uppsala와 Aalborg의 이름을 연결한 도구명이며 임의의 일반 영어 구절로 풀어 쓰는 약어가 아닙니다. 개발기관의 공식 소개에서 두 대학의 공동 개발과 시간 오토마타 기반 도구라는 설명을 확인하고 최초 등장에 명칭과 기능을 덧붙였습니다. 관련 연구에서는 공식 질의 문서와 기존 튜토리얼을 인용했습니다. 도구가 시간 오토마타를 지원하는 사실과 본 중심 모델에 정량 시간 시계가 없는 사실도 구분했습니다.

\end{authorresponse}

\begin{manuscriptrevision}{\mloc{sec:intro}; \mloc{sec:related}; \mloc{sec:queries}; 아래 발췌: 1쪽}

본문에 다음 내용을 반영했습니다.

\begin{quote}

UPPAAL은 Uppsala 대학과 Aalborg 대학이 공동 개발한 시간 오토마타 기반 모델 검사 도구의 이름이며, 두 대학 이름을 연결한 명칭이다.

\end{quote}

\end{manuscriptrevision}

\end{reviewitem}

\clearpage

\begin{reviewitem}{R2-4. 연구 질문과 기여의 연결}

\begin{reviewercomment}

서론의 연구 질문은 본 논문의 연구 동기(motivation)로 보인다. 그렇다면 연구 질문에 대한 답이 본 논문의 기여에 해당하는지를 명확하게 기술해야 한다.

\end{reviewercomment}

\begin{authorresponse}

동기와 검사 가능한 질문을 나누었습니다. 동기는 CoC 작성 중 조건이 누락되거나 앞선 요구와 충돌할 수 있다는 문제입니다. 질문은 의무와 해제 증거를 연결했을 때 모순/근거 부족을 구분하고 최초 사건·의무·사유를 추적할 수 있는가입니다.

질문에 대응하는 기여를 조건별 표현, 상태·속성 실행과 추적, 최초 결과/재시험/자연 한계의 분리로 정리했습니다. 저장된 144건에서 위치와 의무·사유를 대조한 결과가 이에 연결되며, 자연 오류 성능을 확인한 답은 아닙니다.

\end{authorresponse}

\begin{manuscriptrevision}{\mloc{sec:intro}; \mloc{tab:worked-example}; \mloc{tab:retest}; \mloc{sec:conclusion}; 아래 발췌: 1쪽}

본문에 다음 내용을 반영했습니다.

\begin{quote}

첫째, 원문 구절--의무 ID--해제 조건--후속 증거--행동을 연결하고 모순과 근거 부족을 구분하는 표현을 제시한다. 둘째, 조건별 상태와 속성을 UPPAAL로 실행하고 최초 문제 위치와 사유를 추적하는 절차를 구체화한다.

\end{quote}

\end{manuscriptrevision}

\end{reviewitem}

\clearpage

\begin{reviewitem}{R2-5. 서론 구조와 간결성}

\begin{reviewercomment}

서론에서 연구 질문에 대한 구체적인 해결 방법론이 구조적으로 설명되지 않은 것으로 보인다. 전체적으로 장황하므로 간결하고 명확하게 다시 기술할 필요가 있다.

\end{reviewercomment}

\begin{authorresponse}

서론을 작성 오류 가능성→보행자/녹색 신호 예시→조건별 구조화와 UPPAAL 검사→연구 질문→기여와 주장 범위 순으로 재구성했습니다. 전체 흐름 그림에는 자동 텍스트와 작성 프레임의 별도 입력 경로, 검사기와 문제 추적 출력을 표시했습니다.

예시에서 A 미해제 확인은 모순, A 해제 미상은 근거 부족, A 해제 확인은 해당 의무와 양립이라는 구분을 먼저 설명합니다. 상세 알고리즘·입력 사례·질의는 방법 절에 배치했습니다.

\end{authorresponse}

\begin{manuscriptrevision}{\mloc{sec:intro}; \mloc{fig:overview-pipeline}; \mloc{sec:method}; \mloc{tab:worked-example}; 아래 발췌: 1쪽}

본문에 다음 내용을 반영했습니다.

\begin{quote}

녹색 신호는 출발의 새로운 근거이지만 보행자가 통과했다는 증거는 아니다. 통과 여부가 제시되지 않았다면 확정 모순이 아니라 정지 요구의 해제 근거가 부족한 경우이다.

\end{quote}

\end{manuscriptrevision}

\end{reviewitem}

\clearpage

\begin{reviewitem}{R2-6. 관련 연구의 기술 비교}

\begin{reviewercomment}

2장 「관련 연구」에는 관련 연구 동향이나 기존 연구 결과의 문제점, 그리고 이를 보완할 내용이 제시되어야 하나, 현재는 충분히 서술되지 않았다. 이에 대한 보완이 필요하다.

\end{reviewercomment}

\begin{authorresponse}

관련 연구를 원문·공식 자료로 확인한 기술 비교로 보완했습니다. 설명 충실성 연구(Is VLA Reasoning Faithful?)와 VLADriveBench의 관측·개입 검증도 저장 주석의 조건 연결 검사와 비교했습니다. DriveLM/DriveVLM의 설명·계획 생성, runtime verification의 관측 경로 검사, LTL/TLTL의 유한 접두 경로 세 값 의미, UPPAAL의 상태 속성 검사와 Scenic/VerifAI의 상황 생성·분석을 입력·표현·출력으로 비교했습니다.

이전 상태 기억과 세 값·최초 위반이 선행 기술임을 명시했습니다. 기존 방법이 CoC 검사를 할 수 없다고 단정하지 않고, 본 논문이 원문 구절과 의무·조건·증거·문제 기록의 연결을 실증한다는 범위로 보완 방향을 설명했습니다. 문헌 비교를 직접 수행하지 않은 성능 우위로 바꾸지 않았습니다.

\end{authorresponse}

\begin{manuscriptrevision}{\mloc{sec:related}; \mloc{tab:related-position}; \mloc{sec:discussion}; 아래 발췌: 2쪽}

본문에 다음 내용을 반영했습니다.

\begin{quote}

따라서 이전 상태의 기억, 세 값의 판정과 최초 위반 추적 자체는 기존 기술이다. 본 논문의 미상은 명시된 해제 증거의 부족을 뜻하며, 미래 연장 전체에 대한 LTL 세 값 의미론을 구현했다는 뜻은 아니다.

\end{quote}

\end{manuscriptrevision}

\end{reviewitem}

\clearpage

\begin{reviewitem}{R2-7. 기존 연구 대비 기여 범위}

\begin{reviewercomment}

관련 연구 및 기존 연구와 비교하여 본 논문 연구 결과의 독창성이 명확하게 제시되지 않은 것으로 보인다. 이에 대한 보완이 필요하다.

\end{reviewercomment}

\begin{authorresponse}

기여를 상태 기억이나 모델 검사 기법 자체의 발명으로 표현하지 않았습니다. 본 논문은 CoC 원문 구절에서 의무 ID·해제 조건·조건별 증거·행동 전환·최초 문제 기록을 연결해 검토 가능한 입력과 출력으로 만드는 데 초점을 둡니다. 보행자 세 변형의 끝까지 연결된 예시와 실제 질의/출력 대조로 구체화했습니다.

단순 FSM도 성공한 사실을 유지하여 UPPAAL 정확도 우위, 최초·최고 또는 기존 기술로 불가능하다는 주장은 하지 않았습니다. 현재 근거는 작성 시나리오에서의 방법 실증이며 일반적인 독창성·성능 주장은 제한했습니다.

\end{authorresponse}

\begin{manuscriptrevision}{\mloc{sec:related}; \mloc{tab:related-position}; \mloc{tab:worked-example}; \mloc{sec:queries}; \mloc{sec:discussion}; 아래 발췌: 2쪽}

본문에 다음 내용을 반영했습니다.

\begin{quote}

기술적 초점은 원문 구절과 의무ㆍ해제 조건ㆍ증거ㆍ문제 기록을 이어 검토 가능하게 만드는 데 있다. 기존 형식 검증 기술을 새로 발명했다고 주장하지 않으며, 다른 방법이 이러한 연결을 할 수 없다고 단정하지 않는다.

\end{quote}

\end{manuscriptrevision}

\end{reviewitem}

\section*{맺음말}

심사 의견에 따라 의미 변환과 상태·질의 설명, 최초 실패와 재시험, 자료 출처 및 주장 범위를 구체화했습니다. 자연 오류 정확도와 독립 시험, 사람 신뢰도 개선이 미완료라는 한계는 유지했습니다. 이 답변서는 실제 개정 내용과 보존 근거에 대응하며 모든 요구를 독립 성능 검증으로 해결했다고 주장하지 않습니다.

\end{document}
